% Generated draft: requires mathematical and source review.
\documentclass[11pt]{article}
\usepackage[margin=25mm]{geometry}
\usepackage{fontspec}
\setmainfont{Latin Modern Roman}
\usepackage{amsmath,amssymb,mathtools,unicode-math}
\setmathfont{Latin Modern Math}
\usepackage{enumitem,needspace,xurl,hyperref}
\setlength{\parindent}{0pt}
\setlength{\parskip}{5pt}
\setlength{\emergencystretch}{3em}
\newcommand{\N}{\mathbb{N}}
\newcommand{\Z}{\mathbb{Z}}
\newcommand{\Q}{\mathbb{Q}}
\newcommand{\R}{\mathbb{R}}
\newcommand{\C}{\mathbb{C}}
\newcommand{\F}{\mathbb{F}}
\newcommand{\A}{\mathbb{A}}
\newcommand{\PP}{\mathbb{P}}
\newcommand{\E}{\mathbb{E}}
\newcommand{\eps}{\varepsilon}
\newcommand{\dd}{\,\mathrm{d}}
\begin{document}
\section*{UnsolvedMath ID: 10400208}
\textbf{Source problem code:} catalogue entry
\subsection{Definitions and mathematical statement}
A Legendrian knot in a contact \(3\)-manifold is an embedded circle everywhere tangent to its contact plane field. R. Benedetti's conjecture, retaining the source's unresolved notation, is
\[
 \text{for every }W\text{ and every admissible }(pL)\text{-class }\alpha_1,
 \qquad f_1^*\text{ is an isomorphism}.
\]
For whatever objects \(A_{W,\alpha_1}\) and \(B_{W,\alpha_1}\) are the intended domain and codomain, this requires \(f_1^*:A_{W,\alpha_1}\to B_{W,\alpha_1}\) to have an inverse in their specified category. The excerpt omits the preceding definition of \(W\), the \((pL)\)-classes, the map \(f_1\), and its induced map \(f_1^*\); neither its domain nor its codomain can be supplied without inventing context.

The source states the intended consequence that finite-type invariants of Legendrian knots should be insensitive to geometric or rigid contact-structure properties such as tightness, and refers to [50] for more discussion and related questions. This is retained as the reported implication of the missing map statement, not independently strengthened into an assertion about every contact structure or every knot invariant.

A subsequent heading, ``Knots and finite groups,'' begins separate background: knot groups are residually finite, meaning that each nonidentity element survives under some homomorphism to a finite group. Together with the stated criterion that a knot is trivial exactly when its longitude is trivial in the knot group, the source concludes that each nontrivial knot has a finite-group homomorphism detecting its longitude. These are contextual statements from the following subsection, not extra hypotheses or conclusions of Benedetti's conjecture; their usual interpretation concerns knots in \(S^3\) with preferred longitude.
\subsection{Short English statement}
For every source-admissible \(W\) and \((pL)\)-class, Benedetti conjectures that the missing induced map \(f_1^*\) is an isomorphism, with the stated implication that finite-type Legendrian invariants do not detect contact tightness.
\subsection{Sources}
\begin{itemize}
\item R. Benedetti, attributed Conjecture 12.2 in T. Ohtsuki, \emph{Problems on invariants of knots and 3-manifolds} (2002), supplied PDF page 159, together with the following subsection 12.2 background. \url{https://msp.org/gtm/2002/04/gtm-2002-04-024s.pdf}. Recorded access 2026-07-29; extraction pdf-font-structure.
\item The conjecture refers to [50] for detailed discussion and related questions. The bibliography entry and preceding definitions are absent. The following background names Dehn's lemma and the loop theorem without additional references.
\item Local UnsolvedMath registry, \texttt{data/unsolvedmath/problems.json}, permanent ID \(10400208\), exact source code \texttt{catalogue entry}; research classification OPEN-TRIAGE.
\end{itemize}
\end{document}
